\documentclass{amsart}
\input{C:/Users/jzchr/MathDocuments/preamble_general.tex}

\usepackage{enumitem}

\DeclareMathOperator{\PA}{\mathsf{PA}}
\DeclareMathOperator{\ACA}{\mathsf{ACA}}
\DeclareMathOperator{\ZFC}{\mathsf{ZFC}}

\title{The Paris-Harrington Theorem}
\author{Jalen Chrysos}

\begin{document}
	
	\maketitle
	
	\section{The Paris-Harrington Statement} 
	
	Some terminology and notation:
	\begin{itemize}[leftmargin=10pt]
		\item $[X]^n$ denotes the set of ordered $n$-tuples in $X$, and $[d]:= \{i:i<d\}$.
		\item Given $f:[X]^n\to Z$, a subset $Y\subseteq X$ is $f$-\textit{homogeneous} if $f$ is constant on $[Y]^n$.
		\item A function $f:[X]^n\to Z$ is \textit{regressive} if $f(\bar{x})\leq \min(\bar{x})$ for all $\bar{x}\in [X]^n$.
		\item Given $f:[X]^n\to Z$, $Y\subseteq X$ is $f$-\textit{min-homogeneous} if $f(\bar{x})$ only depends on $\min(\bar{x})$.
		\item $X\to (m)^n_k$ means ``For $f:[X]^n\to k$, there is an $f$-homogeneous set $Y\subset X$ with $|Y|\geq m$.''
		\item $X\to (m)^n_{\text{reg}}$ means ``For regressive $f:[X]^n\to \N$, there is an $f$-min-homogeneous set $Y\subset X$ with $|Y|\geq m$.''\\
	\end{itemize}
	
	\noindent \textbf{Finite Ramsey Statement (FR)}: For all $k,m,n$, there exists some $d$ such that for any function $c:[d]^n\to k$, there is a $c$-homogeneous set $Y\subset [d]$ with $|Y|\geq m$.\\
	
	\noindent \textbf{Paris-Harrington Statement (PH)}: For all $k,m,n$, there exists $d$ such that for any function $c:[d]^n\to k$, there is a $c$-homogeneous set $Y\subset [d]$ with $|Y|\geq m$ \textit{and} $|Y|\geq \min(Y)$.\\
	
	There is also a statement equivalent to PH, which we denote PH$\ast$:\\
	
	\noindent \textbf{Paris-Harrington Equivalent (PH$\ast$)}: For all $m,n$, there exists some $d$ for which $[d]\to (m)^n_{\text{reg}}$.\\
	
	The Finite Ramsey Statement is provable in Peano Arithmetic ($\PA$), but Paris and Harrington showed that the Paris-Harrington Statement is not, though it \textit{is} provable in stronger theories ($\ZFC$, for example), so it is still considered true. 
	
	Their proof showed that in $\PA$, PH implied the consistency of $\PA$, violating G\"odel's Second Incompleteness Theorem. So either $\PA$ itself was inconsistent, or it could not prove PH. 
	
	Later, Kanamori and McAloon gave another proof using model theory. They constructed a model of $\PA$ in which PH was false, thereby showing that $\PA$ could not prove PH. This short write-up will follow their approach.
	
	\newpage 
	
	\section{Outline of the Argument}
	
	Kanamori and McAloon's approach constructs a model of $\PA$ by using \textit{diagonal indiscernibles}:\\
	
	\noindent \textbf{Definition (Diagonal Indiscernibles)}: In a structure $\M$, given a set of sentences $\Delta$, a sequence $I\subset \M$ is \textit{diagonal-indiscernible} relative to $\Delta$ if for any $\phi(a;\bar{x})\in \Delta$ (where $a$ is a parameter in $\M$) and elements $b<\bar{y},\bar{z}\in I$, we have $\M\models \phi(a;\bar{y}) \iff \phi(a;\bar{z})$ for all $a<b$.\\
	
	The proof requires two lemmas. The first shows the existence of certain diagonal indiscernible sequences, and the second shows how such sequences can be used to create models of $\PA$.\\
	
	\noindent \textbf{Lemma 1}: For any $m,n,\ell$, for any sequence of formulas $\Delta=\{\phi_{\a}(a;\bar{x})\}_{\a<\ell}$ where $|\bar{x}|\leq n$, there is a set $|I|\geq m$ that is diagonal-indiscernible relative to $\Delta$. In particular, such a sequence exists within $[c,d]$ when $[c,d]\to (h)_{\text{reg}}^{2n+1}$, where $h\to (m+n)^{2n+1}_{\ell+1}$.\\
	
	\noindent \textbf{Lemma 2}: Given a model $\M\models \PA$ with an infinite set $I$ that is diagonal-indiscernible relative to all $\Delta_0$ sentences, the substructure
	$$
	\NN(I) := \bigcup_{i\in I} \{x\in \MM : x<i\}
	$$
	is also a model of $\PA$.\\
	
	With these two Lemmas, we have:\\
	
	\noindent \textbf{Lemma 3}: Let $\M\models \PA$ and $c\in \M$ be non-standard. If $d\in \M$ is such that $[c,d]\to (2c)^c_{\text{reg}}$, then there is an initial segment of $d$ including $c$ which models $\PA$.
	\begin{proof}
		Let $F(c,d;r)$ represent the following statement: setting $m=n=\ell=r$ as in Lemma 1, diagonal indiscernibles $I$ can always be found in $[c,d]$. $F(c,d;r)$ holds for any standard $r$ as follows: first, the minimum $h$ for which $h\to (m+n)^{2n+1}_{\ell+1}$ is standard by FR. Next, for any standard $p,q$,
		$$\Big(A\to (2c)^c_{\text{reg}}\Big) \implies \Big(A\to (p)^{q}_{\text{reg}}\Big).$$
		To prove this, consider a regressive function $f:A^q\to \N$, one can extend it to a function $\bar{f}:A^{c}\to \N$ via $\bar{f}(\bar{x})=f(\bar{x}|_q)$. $A$ has a $\bar{f}$-min-homogeneous set of size $2c$ which we can call $Y$, so let its first $c$ elements be $Y'$. Then $Y'$ is $f$-min-homogeneous. Thus, setting $A=[c,d]$, $p=h$, and $q=2n+1$, we see that the condition in Lemma 1 is fulfilled.
		
		Moreover, $F(c,d;r^*)$ must also hold for some non-standard $r^*$; if not, then $F(c,d;r)$ exactly defines the standard numbers, so by internal induction in $\PA$, all of $\M$ is standard, a contradiction\footnote{This type of argument is sometimes called Robinson's ``overspill'' principle.}.
		
		Indexing the formulas of $\Delta_0$ by $r^*$ (a countable initial segment of $r^*$ will do), we have an infinite diagonal-indiscernible sequence $I\subset [c,d]$ relative to $\Delta_0$. By Lemma 2, this determines a model $\NN(I)\models \PA$, and $\NN(I)$ contains $c$ but not $d$.
	\end{proof}\\
	
	From Lemma 3, the theorem follows:\\
	
	\noindent \textbf{Theorem (Paris-Harrington)}: PH is not provable in $\PA$.
	\begin{proof}
		Let $\M\models \PA$ be a nonstandard model, and choose some non-standard element $c$. Assuming PH holds in $\M$, let $d$ be minimal such that $[c,d]\to (2c)^c_{\text{reg}}$, so $d$ is also non-standard. Let $\NN$ be the model of $\PA$ given by Lemma 3, so that $c\in\NN$ but $d\not\in \NN$. I claim PH does not hold in $\NN$.
		
		Indeed, $\NN$ does not contain $d$, so for all $d'\in \NN$, $\M$ believes $[c,d']\not\to (2c)^c_{\text{reg}}$; that is, there is a regressive function $f:[c,d']^c\to \M$ with no $f$-min-homogeneous set $|Y|\geq 2c$ in $\M$, and hence none in $\NN$ either. A priori, it may be that such an obstruction exists only in $\M$ and not in $\NN$, but because $f$ is \textit{regressive}\footnote{This seems to be a place where working with PH$\ast$ rather than PH is necessary.}, the codomain of $f$ lies in $\NN$, so $\NN\models [c,d']\not\to (2c)^c_{\text{reg}}$. Thus, PH$\ast$ and hence PH both fail in $\NN$.
	\end{proof}\\
	
	\newpage 
	
	\section{Proving Lemmas 1 \& 2}
	
	\noindent \textbf{Lemma 1}: For any $m,n,\ell$, for any sequence of formulas $\Delta=\{\phi_{\a}(a;\bar{x})\}_{\a<\ell}$ where $|\bar{x}|\leq n$, there is a set $|I|\geq m$ that is diagonal-indiscernible relative to $\Delta$. In particular, such a sequence exists within $[c,d]$ when $[c,d]\to (h)_{\text{reg}}^{2n+1}$, where $h\to (m+n)^{2n+1}_{\ell+1}$.\\
	\begin{proof} 
		First, let's see how Ramsey's theorem enables us to find indiscernible sequences in general, and then we'll make a slightly stronger argument to get sequences of diagonal indiscernibles.
		
		Suppose the formulas $\phi_{\a}$ take no parameters, i.e. $k=0$. We define a function $g:[\N]^{2n}\to [\ell+1]$. For a $2n$-tuple $\{x_1<\cdots<x_n<y_1<\cdots<y_n\}$, $g$ is defined
		$$
		g(\bar{x},\bar{y}) := \begin{cases}
			\ell + 1 & \text{if } \phi_{\a}(\bar{x})\iff \phi_{\a}(\bar{y}) \; \forall \a<\ell,\\
			\min \{\a<\ell : \phi_{\a}(\bar{x}) \neq \phi_{\a}(\bar{y})\} & \text{else}.
		\end{cases}
		$$
		That is, $g$ detects the first failure of indiscernibility. If $X$ is $g$-homogeneous and $g(X)=\ell + 1$, this is not quite enough to guarantee indiscernibility because $g$ only compares tuples where $\bar{x}$ is entirely below $\bar{y}$, and for indiscernibility one must consider tuples that are interleaved together. However, if we take such an $X$ and remove the largest $n$ elements $\bar{t}$, the resulting $I$ \textit{is} indiscernible as $\phi_{\a}(\bar{x})=\phi_{\a}(\bar{t})=\phi_{\a}(\bar{y})$ for any $\bar{x},\bar{y}$. Thus, it suffices to find such an $X$ with $|X|\geq m+n$.
		
		Using Ramsey's Theorem, it's possible to find $g$-homogeneous sets of any size, and if $X$ is $g$-homogeneous with $|X|\geq 3n$ then $g(X)=\ell+1$ is forced: if $g(X)=i\leq \ell$, then there are disjoint $\bar{x},\bar{y},\bar{z}\in X$ so that $\phi_i(\bar{x}),\phi_i(\bar{y}),\phi_i(\bar{z})$ are all mutually unequal, but this is impossible as there are only 2 truth values.\\
		
		Now to extend this to diagonal indiscernibles. The previous argument now faces an issue: it may be that the truth values of $\phi_{i}(a;\bar{x}),\phi_i(a;\bar{y}),\phi_i(a;\bar{z})$ \textit{are} all mutually unequal, but for different choices of the parameter $a$. So we must take the additional step of finding a sufficiently large $H\subset \N$ on which a single $a$ witnesses discernibility. 
		
		Define $g:[\N]^{2n+1}\to [\ell+1]$ similarly to before, so that if $\{b<x_1<\cdots<x_n<y_1<\cdots<y_n\}$, $g(b,\bar{x},\bar{y})$ is the least $\a$ so that $\phi_{\a}(a;\bar{x})\neq \phi_{\a}(a;\bar{y})$ for some $a<b$ if such $\a$ exists, and $\ell+1$ otherwise. Also define $f:[\N]^{2n+1}\to \N$ so that $f(b,\bar{x},\bar{y})$ is the least $a$ where $\phi_{\a}(a;\bar{x})\neq \phi_{\a}(a;\bar{y})$ for $\a=g(b,\bar{x},\bar{y})$, if $\a\neq \ell+1$, and $b$ otherwise.
		%		$$
		%		f(b,\bar{x},\bar{y}) := \begin{cases}
			%			(b,b,\dots,b) & \text{if } \bar{x},\bar{y} \text{ are d.i. wrt } \{\phi_{\a}\}_{\a<\ell}\\
			%			\min\{\bar{a}<b : \phi_{\a}(\bar{a};\bar{x}) \neq \phi_{\a}(\bar{a};\bar{y})\} & \text{else}.
			%		\end{cases}
		%		$$
		Note that $f$ is regressive. This implies that there is an arbitrarily large $H\subseteq \N$ that is $f$-min-homogeneous.
		
		Now we proceed as in the first case, choosing a $g$-homogeneous $|X|\geq m+n,3n+1$ from within $H$ so that we have some disjoint $b < \bar{x} < \bar{y}< \bar{z}$. Let $a=f(b,\bar{x},\bar{y})=f(b,\bar{x},\bar{z})=f(b,\bar{y},\bar{z})$ (they are equal by min-homogeneity of $H\supset X$), and $i=g(b,\bar{x},\bar{y})$. Then $\phi_i(a;\bar{x}),\phi_i(a;\bar{y}),\phi_i(a;\bar{z})$ are all pairwise distinct, but again this is impossible.\\
		
		Finally, in which intervals $[c,d]$ is such an $I$ guaranteed to exist? To find a $g$-homogeneous $|X|\geq m+n$ within $H$ requires $H\to (m+n)^{2n+1}_{\ell+1}$, so let $h$ be the minimum size of $H$ for which this holds. To guarantee that an $f$-min-homogeneous set of this size exists in $[c,d]$, we require $[c,d]\to (h)_{\text{reg}}^{2n+1}$. 
%		
%		Note that $\big(d\to (2c)^c_{\text{reg}}\big)\implies \big(d\to c^p_{\text{reg}}\big)$ for $p<c$, as follows: for any regressive $f:[d]^p\to [d]$, we can extend it to a regressive function $\bar{f}:[d]^c\to [d]$ by
%		$$
%		\bar{f}(x_1,\dots,x_c) := f(x_1,\dots,x_p).
%		$$
%		Because $d\to (2c)^c_{\text{reg}}$, there is some $X\subset [d]$ with $|X|=2c$ that is $\bar{f}$-min-homogeneous. Let $Y$ be the first $c$ elements of $X$. Then $Y$ is $f$-min-homogeneous: any two $p$-tuples in $Y$ with the same minimum can be completed to $c$-tuples in $X$, which therefore have the same $\bar{f}$ value, and hence the same $f$ value as well.
	\end{proof}\\
	
%	\noindent \textbf{Lemma 1}: For any $m,n,k$, given a sequence of sentences $\{\phi_{\a}(\bar{a};\bar{x})\}_{\a<\ell}$ where $|\bar{a}|=k,|\bar{x}|=n$, there is a set $|I|\geq m$ that is diagonal-indiscernible relative to $\{\phi_{\a}\}$. In particular, if $[c,d]\to (2c)^{c}_{\text{reg}}$, then $I$ can be chosen from within $(c,d)$.
%	\begin{proof}
%		First, let's see how Ramsey's theorem enables us to find indiscernible sequences in general, and then we'll make a slightly stronger argument to get sequences of diagonal indiscernibles.
%		
%		Suppose the formulas $\phi_{\a}$ take no parameters, i.e. $k=0$. We define a function $g:[\N]^{2n}\to [\ell+1]$. For a $2n$-tuple $\{x_1<\cdots<x_n<y_1<\cdots<y_n\}$, $g$ is defined
%		$$
%		g(\bar{x},\bar{y}) := \begin{cases}
%			\ell + 1 & \text{if } \phi_{\a}(\bar{x})\iff \phi_{\a}(\bar{y}) \; \forall \a<\ell,\\
%			\min \{\a<\ell : \phi_{\a}(\bar{x}) \neq \phi_{\a}(\bar{y})\} & \text{else}.
%		\end{cases}
%		$$
%		That is, $g$ detects the first failure of indiscernibility. If $X$ is $g$-homogeneous and $g(X)=\ell + 1$, this is not quite enough to guarantee indiscernibility because $g$ only compares tuples where $\bar{x}$ is entirely below $\bar{y}$, and for indiscernibility one must consider tuples that are interleaved together. However, if we take such an $X$ and remove the largest $n$ elements $\bar{t}$, the resulting $I$ \textit{is} indiscernible as $\phi_{\a}(\bar{x})=\phi_{\a}(\bar{t})=\phi_{\a}(\bar{y})$ for any $\bar{x},\bar{y}$. Thus, it suffices to find such an $X$ with $|X|\geq m+n$.
%		
%		Using Ramsey's Theorem, it's possible to find $g$-homogeneous sets of any size, and if $X$ is $g$-homogeneous with $|I|\geq 3n$ then $g(X)=\ell+1$ is forced: if $g(X)=i\leq \ell$, then there are disjoint $\bar{x},\bar{y},\bar{z}\in X$ so that $\phi_i(\bar{x}),\phi_i(\bar{y}),\phi_i(\bar{z})$ are all mutually unequal, but this is impossible as there are only 2 truth values.\\
%		
%		Now to extend this to diagonal indiscernibles. The previous argument now faces an issue: it may be that the truth values of $\phi_{i}(\bar{a};\bar{x}),\phi_i(\bar{a};\bar{y}),\phi_i(\bar{a};\bar{z})$ \textit{are} all mutually unequal, but for different choices of the parameters $\bar{a}$. So we must take the additional step of finding a sufficiently large $H\subset \N$ on which a single $\bar{a}$ witnesses discernibility. 
%		
%		Define $g:[\N]^{2n+1}\to [\ell+1]$ similarly to before, so that if $\{b<x_1<\cdots<x_n<y_1<\cdots<y_n\}$, $g(b,\bar{x},\bar{y})$ is the least $\a$ so that $\phi_{\a}(\bar{a};\bar{x})\neq \phi_{\a}(\bar{a};\bar{y})$ for some $\bar{a}<b$ if such $\a$ exists, and $\ell+1$ otherwise.
%		
%		Also define $f_1,\dots,f_k:[\N]^{2n+1}\to \N$ so that 
%		$$
%		\{f_1(b,\bar{x},\bar{y}),f_2(b,\bar{x},\bar{y}),\dots,f_k(b,\bar{x},\bar{y})\}
%		$$
%		is some $\bar{a}$ where $\phi_{\a}(\bar{a};\bar{x})\neq \phi_{\a}(\bar{a};\bar{y})$ for $\a=g(b,\bar{x},\bar{y})$, if $\a\neq \ell+1$, and $(b,\dots,b)$ otherwise.
%%		$$
%%		f(b,\bar{x},\bar{y}) := \begin{cases}
%%			(b,b,\dots,b) & \text{if } \bar{x},\bar{y} \text{ are d.i. wrt } \{\phi_{\a}\}_{\a<\ell}\\
%%			\min\{\bar{a}<b : \phi_{\a}(\bar{a};\bar{x}) \neq \phi_{\a}(\bar{a};\bar{y})\} & \text{else}.
%%		\end{cases}
%%		$$
%		Note that each $f_j$ is regressive. This implies that there is an infinite $H\subseteq \N$ that is min-homogeneous for all of them.
%		
%		Now we proceed as in the first case, choosing a $g$-homogeneous $|X|\geq m+n,3n+1$ from within $H$ so that we have some disjoint $b < \bar{x} < \bar{y}< \bar{z}$. Let $\bar{a}=f(b,\bar{x},\bar{y})=f(b,\bar{x},\bar{z})=f(b,\bar{y},\bar{z})$ (they are equal by min-homogeneity of $H\supset X$), and $i=g(b,\bar{x},\bar{y})$. Then $\phi_i(\bar{a};\bar{x}),\phi_i(\bar{a};\bar{y}),\phi_i(\bar{a};\bar{z})$ are all distinct, but again this is impossible.\\
%		
%		Now there's a question which will be important later: how small can we keep $\max(I)$? Suppose we want to choose $I$ from within an initial segment $[d]$. To find a $g$-homogeneous $|X|\geq m+n$ within $H$ requires $H\to (m+n)^{2n+1}_{\ell+1}$, so let $h$ be the minimum size of $H$ for which this holds. To guarantee that an $f_j$-min-homogeneous set of this size exists, we require $d\to (h)_{\text{reg}}^{2n+1}$.
%	\end{proof}\\
	\newpage 
	
	\noindent \textbf{Lemma 2}: Given a model $\M\models \PA$ with an infinite set $I$ that is diagonal-indiscernible relative to all $\Delta_0$ sentences, the substructure
	$$
	\NN(I) := \bigcup_{i\in I} \{x\in \MM : x<i\}
	$$
	is also a model of $\PA$.
	\begin{proof}
		Let $I=\{x_1<x_2<\cdots\}$. Being diagonal-indiscernible over all of $\Delta_0$ means roughly that from the ``viewpoint'' of $x_i$, there is no bounded-quantifier way to distinguish $x_{i+1}$ from $x_{i+2}$, and so on. In particular, $\NN(I)$ is closed under addition and multiplication as follows:\\
		
		Suppose $i<j<k$ and $x_i+x_j>x_k$. Then let $a:=x_k-x_j$, and note $a<x_i$, so diagonal indiscernibility forces any two members of $I$ to have difference $a$, but this is impossible. This shows that $\NN(I)$ is closed under addition.
		
		Similarly, if $x_i\cdot x_j > x_k$ then letting $a=\lfloor x_k/x_j\rfloor$ so that $a\cdot x_j < x_k \leq (a+1)x_j$, we have $a<x_i$. Diagonal indiscernibility forces any two members of $I$ to have ratio in $[a,a+1)$, which is impossible (using the fact that $a\geq 2$). This shows that $\NN(I)$ is closed under multiplication.\\
		
		It remains to show that $\NN(I)$ satisfies induction. To do this, we first show that formulas in $\NN(I)$ have corresponding equivalent formulas in $\M$. Consider an arbitrary sentence 
		$$\phi(a) = (\exists v_1)(\forall v_2) \cdots (\exists v_n): \psi(a,\bar{v})$$
		where $\psi\in \Delta_0$. For any $a\in \NN(I)$ with say $a<x_i$, the sentence
		$$
		\phi_0(a;\bar{y}):= \; (\exists v_1 < y_1)(\forall v_2 < y_2)\cdots (\exists v_n<y_n): \psi(a,\bar{v})
		$$
		is $\Delta_0$, so for all $\bar{y}>a$ its truth value in $\NN(I)$ is constant, by diagonal indiscernibility of $I$. And the truth value of $\phi_0(a;\bar{y})$ for sufficiently large $\bar{y}$ is equal to that of $\phi(a)$ in $\NN(I)$, as $I$ is unbounded in $\NN(I)$; the bound $\bar{y}$ will eventually be large enough to see any given witness $\bar{v}$. Thus, 
		$$
		\NN(I)\models \phi(a) \iff \NN(I) \models \phi_0(a;\bar{y}) \;\, \text{for }\bar{y}>a.
		$$ Moreover, because $\NN(I)$ is an initial segment of $\M$, they agree on $\Delta_0$ formulas, so in fact
		$$
		\NN(I)\models \phi(a) \iff \M \models \phi_0(a;\bar{y}) \;\, \text{for }\bar{y}>a.
		$$
		Induction is equivalent to the statement that for any $\phi(a)$, there is a least $a$ where $\phi(a)$ holds. Because $\MM\models \PA$ and thus has induction, there is a least $a$ such that $\phi_0(a;\bar{y})$ holds, so the same is true for $\NN(I)$.
	\end{proof}
	
\end{document}